\section{Experimental setup}
\label{sec:experimental-setup}

\subsection{Research questions and evaluation design}

The empirical study addresses the five prespecified questions in
Table~\ref{tab:research-questions}.  Each question is linked to an estimand, a
comparison set, and uncertainty summaries that respect dependence among states
from the same optimization trajectory and function family.

\begin{table*}[t]
  \caption{Research questions and corresponding evaluation criteria.}
  \label{tab:research-questions}
  \centering
  \small
  \begin{tabularx}{\textwidth}{@{}c Y Y@{}}
    \toprule
    RQ & Question & Evaluation criteria \\
    \midrule
    RQ1 & Statewise value of the primary query & Utility distribution and decomposition; proportion with $U_q\le0$; family-balanced effects and intervals. \\
    RQ2 & Information beyond elapsed budget & Nested training-family OOF selection; paired B3--T0 effects; frozen BBOB-validation evaluation. \\
    RQ3 & End-to-end policy trade-off & Call quality, final performance, FE use, runtime, and paired comparisons with all required policies. \\
    RQ4 & Transport of the frozen procedure & Separate held-out BBOB, CEC2017, CEC2022, and engineering estimates, with coverage and query failures. \\
    RQ5 & Behavioral associations with query value & T0--B3 increments, model-parameter stability, maturity--utility association, and prespecified strata. \\
    \bottomrule
  \end{tabularx}
\end{table*}

\subsection{Internal benchmark and function-family partition}

The internal benchmark is the noiseless BBOB suite.  The training set comprises
functions 1, 2, 3, 4, 6, 7, 8, 10, 11, 12, 15, 16, 17, 18, 20, 21, 22, and 23;
functions 5, 9, 13, 14, 19, and 24 form the held-out validation set.  Functions,
rather than randomly selected transformed instances, define the grouping
variable.  This partition addresses evidence that instance-level splits can
overstate algorithm-selection generalization
\cite{cenikjLandscapeFeaturesSingleobjective2025}.  Both sets use dimensions
$d\in\{10,20,40\}$, instances $1,2,3$, optimizer seeds $1,\ldots,30$, population
size 40, and total budget $B=1000d$.  The portfolio contains DE, PSO, CMA-ES,
and SHADE, giving 54 training and 18 validation function--dimension
combinations.

The single best solver (SBS) is estimated exclusively from terminal outcomes at
$FE=B$.  For each problem--algorithm pair, final best fitness is first averaged
over seeds.  The four algorithms are then ranked within each problem, with
smaller mean fitness preferred and average ranks assigned to ties.  These ranks
are averaged over the training problems; any remaining tie is resolved in the
prespecified order DE, PSO, CMA-ES, SHADE.  Intermediate decision states,
including the last state at or before approximately $0.60B$, do not contribute
to the SBS estimate.

BBOB validation is a held-out internal evaluation set.  It remains untouched
during query selection, preprocessing, model-family selection, input-group
comparison, threshold fitting, and score-neighborhood definition.

\subsection{External benchmarks}

External evaluation is organized as non-pooled, suite-specific analyses.  The
CEC2017 analysis comprises functions 1--29, dimensions 10, 30, and 50, seeds
$1,\ldots,30$, population size 40, and budget $B=1000d$.  The four-algorithm
portfolio and dynamic decision-opportunity protocol are unchanged, and the
controller, downstream selector, imputation values, and threshold are held
fixed at their BBOB-training estimates.

CEC2022 and engineering problems form separate components of RQ4.  Their
problem sets, dimensions, bounds, budgets, success targets, and failure rules
are specified independently of policy-outcome inspection.  For every external
suite, query-extraction failures remain in the coverage denominator rather than
being removed from the analysis.  Missing descriptors may be imputed only with
BBOB-training medians.  Held-out BBOB, CEC2017, CEC2022, and engineering
estimates are reported separately; neither partial coverage nor agreement on a
subset of suites is interpreted as unconditional cross-benchmark transfer.

\subsection{Landscape-query configurations}

Table~\ref{tab:query-configurations} defines the three prespecified queries.
Primary inference concerns \texttt{descriptor\_cheap}; the two \texttt{pflacco}
configurations evaluate dependence on representation and acquisition cost.

\begin{table*}[t]
  \caption{Prespecified landscape-descriptor queries.}
  \label{tab:query-configurations}
  \centering
  \small
  \begin{tabularx}{\textwidth}{@{}l l c c X@{}}
    \toprule
    Query & Sample & FE share & Features & Role \\
    \midrule
    \texttt{descriptor\_cheap} & LHS $50d$ & 5\% & 16 & Primary \\
    \texttt{pflacco\_standard} & Same LHS $50d$ & 5\% & 37 & Configuration robustness \\
    \texttt{pflacco\_broad} & Independent LHS $100d$ & 10\% & 52 & Configuration robustness \\
    \bottomrule
  \end{tabularx}
\end{table*}

The cheap and standard queries use the same evaluated $(X,y)$ sample, yielding
a paired comparison of descriptor representations at a common FE cost.  The
broad query uses an independent $100d$ design.  The standard configuration
comprises PCA, nearest-better clustering, dispersion, information-content, and
ELA-distribution groups.  The broad configuration adds ELA level-set and
sample-derived fitness--distance-correlation groups, while excluding the full
quadratic meta-model group, cell mapping, and feature groups that require
additional objective evaluations.  Query sample points are not inserted into
the optimizer population.  Each query defines a separate downstream selector,
utility target, decision model, threshold, baseline comparison, and external
evaluation; observations from different query identities or FE designs are not
pooled.

\subsection{Dynamic decision opportunities}

Offline label construction and policy evaluation use the same dynamic
budget--event sampling scheme.  Complete native updates are monitored when they
cross a grid from $0.20B$ to $0.60B$ in increments of $0.01B$.  The mandatory
budget milestones are 0.20, 0.22, 0.24, 0.26, 0.28, 0.30, 0.34, 0.38, 0.42,
0.46, 0.50, and 0.60.  Five event types may add decision states: resumed strict
improvement; a stagnation span reaching $0.05B$; an absolute medium-window
covariance-effective-rank change reaching $0.20$; unit-cube elite-centroid
displacement divided by $\sqrt{d}$ reaching $0.05$; and relative diversity
recovery reaching $0.10$.  Covariance-effective-rank-change, elite-migration,
and diversity-recovery events rearm when their respective quantities fall below
$0.10$, $0.025$, and $0.05$.  Stagnation rearms after a new strict best-so-far
improvement, whereas resumed improvement rearms when short-window improvement
frequency returns to zero.

Each complete update that crosses one or more monitor points is evaluated once.
If a crossing includes a milestone, the event is associated with that milestone
and does not consume an event-only quota or spacing anchor.  Otherwise, the most
recent crossed monitor point supplies the nominal location.  Each of the early,
middle, and late phases admits at most two event-only states, with at least
$0.02B$ separation in realized FE ratio.  Suppressed crossings still update
event rearming.  This scheme admits 12--18 states per run, each with unit sample
weight.  The realized $FE/B$ represents elapsed budget and is the sole T0
covariate; exact integer FE aligns shared-state outcomes.  Nominal milestones
and native-update counts describe the sampling process but do not add Decision
Model inputs.

\subsection{Matched states and analysis population}

The analysis pairs all alternatives at an identical complete native state.
For each eligible state, it combines permutation-invariant behavior, the shared
$50d$ and independent $100d$ LHS samples, the three descriptor representations,
and the query-budget-specific outcome of every eligible action.  A trajectory
enters an estimand only when both its complete native history and terminal
outcome satisfy the same state definition.  Query-specific selection
references, utilities, decision models, and policies are estimated from these
matched observations under the family partitions described above.

The primary Decision Model population contains states for which both the prefix
and default algorithm equal the training SBS and the no-query path continues
that optimizer.  States generated from the other prefix algorithms are reserved
for cross-prefix robustness, leave-one-prefix-algorithm-out analysis, and tests
of portability across the four search processes.  They are excluded from
primary model fitting, threshold fitting, and primary result aggregation.

\subsection{Comparison policies}

Eight prespecified comparison roles are retained because they answer different
benchmarking questions.
\begin{itemize}
  \item \textbf{Never Query}: continue the complete SBS state natively and incur no query acquisition cost.
  \item \textbf{Always Query}: execute the fixed query at the first accepted decision opportunity and invoke the same downstream selector used by the proposed policy.
  \item \textbf{Random Analysis}: at each shared opportunity before a call, trigger with probability 0.5.  Thirty independent repetitions use explicit integer random streams.
  \item \textbf{Traditional AAS}: execute the fixed query and apply the same statewise downstream selector, without a preceding Decision-before-Feature controller, consistent with the sample--feature--selector structure of landscape-based selection \cite{guoAutomatedAlgorithmSelection2025}.  Under this definition, it induces the same first-opportunity decision as Always Query.
  \item \textbf{SBS}: use the training-derived static default defined above.  In the primary population, the prefix and default are this same SBS, so its native complete-budget path coincides with Never Query; the separate name preserves its role as the standard algorithm-selection reference rather than introducing another online trajectory.
  \item \textbf{VBS}: use the static, per-problem hindsight choice among complete-budget portfolio outcomes.  It is an upper reference rather than a deployable method and is distinct from the statewise best observed action.
  \item \textbf{Time-only Controller}: use the selected Decision Model family and the same training and threshold procedure as the full controller, with input restricted to $\{FE/B\}$.
  \item \textbf{Decision-before-Feature}: use B3 behavior with the selected model and the threshold estimated from full-training OOF predictions.
\end{itemize}
In the primary population, Never Query and SBS yield the same native terminal
outcome, whereas Always Query and Traditional AAS yield the same
first-opportunity query outcome.  RQ3 therefore reports these roles as
\emph{Never Query / SBS} and \emph{Always Query / Traditional AAS}; each shared
outcome enters summaries and inferential comparisons once.
All online policies share the same accepted decision opportunities, total FE
budget, portfolio, query definition, downstream selector, and
population-transfer rule.  State-only, query-only, and combined statewise
selectors are additionally compared as diagnostics of the downstream selector;
they do not replace the required policy baselines.

\subsection{Evaluation measures}

Decision-model family selection uses nested function-family OOF mean decision
utility.  AUROC, average precision, and Spearman correlation are secondary
score-level measures.  Continuous-utility RMSE is reported only for Ridge;
classification probabilities and discriminant scores are not compared with
continuous utility using RMSE.

Policy-level reporting includes query-call rate, mean decision utility, the
proportion of called states for which $U_q>0$, the fraction of positive-utility
states called, and
\begin{equation}
  \overline U_{\mathrm{decision}}(\pi)
  =\frac{1}{n}\sum_{i=1}^{n}d_{\pi,i}U_q(s_i),
  \label{eq:policy-decision-utility}
\end{equation}
where $d_{\pi,i}$ is the call indicator of query-decision policy $\pi$.  A
skipped state contributes zero, so the Never Query role has zero decision
utility by definition.  The coincident SBS reference role and VBS are not
query-decision policies and therefore remain not applicable for this quantity;
mean observed utility conditional on calling is reported separately.
Positive-utility capture is summarized by
\begin{equation}
  \mathrm{utility\ capture}=
  \frac{\sum_i U_{q,i}\,\mathbf{1}[U_{q,i}>0]\,\mathbf{1}[\widehat d_i=1]}
       {\sum_i U_{q,i}\,\mathbf{1}[U_{q,i}>0]}.
  \label{eq:utility-capture}
\end{equation}
The positive-state capture rate and utility capture in
Eq.~\eqref{eq:utility-capture} are reported separately.  Calls with $U_q\le0$
are summarized by their count, their proportion among calls, and their
accumulated utility loss.

End-to-end endpoints comprise final best objective value or error,
target-based success rate and expected running time (ERT), total and query FEs,
wall-clock time, behavior-extraction time, controller-inference time, query
sampling and descriptor-computation time, downstream-selection latency, and
peak memory.  Success targets are suite specific and are fixed independently
of the observed policy contrasts.  Cost--performance reporting retains all
eight prespecified comparison roles but uses six distinct primary-population
outcomes: five online trajectories, including the two slash-labelled
identities, and the VBS hindsight reference.

\subsection{Statistical analysis}

Repeated states from one trajectory are not treated as independent replicates.
Descriptive summaries include state-level micro-averages and macro-averages
balanced across function families.  Results are also stratified by family,
dimension, phase, and transition type.  Support for the RQ1 statement that a
majority of evaluated states do not benefit from the primary query requires the
95\% lower confidence bound for the relevant proportion to exceed 0.5.  Paired
policy comparisons use nonparametric procedures, including Wilcoxon signed-rank
comparisons and Friedman comparisons with Holm-adjusted follow-ups, at a
prospectively declared aggregation level.
For every endpoint, analytically identical roles enter once; duplicated role
labels add no ranks, paired contrasts, or multiplicity counts.

Equivalence is assessed only against a smallest effect of interest, margin, and
primary procedure fixed before policy outcomes are examined.  Confidence
intervals use a prespecified hierarchical resampling unit and replication
count, and ERT uses suite-specific targets fixed independently of policy
outcomes.  Absence of statistical significance is not interpreted as
equivalence or preservation of optimization performance.
